\chapter{Optimisation convexe}\label{ch:qm:convex-optimisation}

À quatre heures moins le quart, un optimiseur doit transformer deux cents prévisions de rendement en un portefeuille long-short avant la clôture~: neutre en dollars et
neutre dans chacun des cinq secteurs, aucune position supérieure à 1.5~\% du capital, exposition brute au plus égale à 100~\%, et au plus 20~\% du capital échangé
contre le carnet de la veille. Il doit renvoyer le meilleur portefeuille que ces limites autorisent, et il doit aussi dire ce que coûte chacune d'elles, car demain
quelqu'un demandera si la limite de rotation en vaut la peine. Le jour où un responsable des risques ajoute une limite qui contredit les autres, il doit affirmer
qu'aucun portefeuille n'existe et en montrer la raison, sous une forme qu'une personne peut vérifier. L'optimisation convexe assure les trois~: une solution, un prix
pour chaque contrainte, et un certificat lorsqu'il n'existe pas de solution. Ce chapitre expose la théorie qui les garantit (convexité, dualité, conditions d'optimalité),
les classes de problèmes dans lesquelles s'inscrit la construction de portefeuille, puis un solveur maison, qu'il applique au portefeuille.

\section{Ensembles, fonctions et problèmes convexes}

\begin{definition}[Ensemble convexe, fonction convexe, problème d'optimisation convexe]\label{def:qm:convex-optimisation:convex}
Un ensemble $C \subseteq \R^n$ est un \emph{ensemble convexe}\index{ensemble convexe} s'il contient le segment reliant deux quelconques de ses points. Une fonction $f$ est une \emph{fonction
convexe}\index{fonction convexe} si $f(\theta x + (1 - \theta)y) \le \theta f(x) + (1 - \theta)f(y)$ pour $\theta \in [0, 1]$. Un \emph{problème d'optimisation convexe}\index{problème d'optimisation convexe}
minimise une fonction convexe sous des contraintes d'inégalité convexes $g_i(x) \le 0$ et des contraintes d'égalité affines
$Ax = b$.
\end{definition}

\begin{proposition}[Les minima locaux sont globaux]\label{prop:qm:convex-optimisation:local}
Pour un problème convexe, tout minimum local est un minimum global, et l'ensemble des minimiseurs est convexe.
\end{proposition}

\begin{proof}
Si $x$ est un minimum local et si $y$ est réalisable avec $f(y) < f(x)$, les points $x + \theta(y - x)$ sont réalisables pour $\theta \in (0, 1]$ et, par convexité, $f(x + \theta(y -
x)) \le f(x) + \theta(f(y) - f(x)) < f(x)$, quel que soit $\theta$ aussi petit soit-il~: contradiction. Si $x$ et $y$ minimisent tous deux, tout point situé entre eux minimise également.
\end{proof}

C'est pourquoi la convexité importe plus que tout algorithme~: un solveur qui s'arrête en un point vérifiant des conditions locales a trouvé la réponse. Les problèmes de portefeuille
sont convexes lorsque l'objectif est le rendement espéré moins un multiple d'une variance construite à partir d'une \omterm{def:qm:covariance-estimation-and-random-matrices:cov}{matrice de covariance} semi-définie positive (chapitre~22), et que les contraintes sont
des limites linéaires, des valeurs absolues bornées supérieurement (exposition brute, rotation) et des normes bornées supérieurement (risque et tracking error).

\section{Dualité}

\begin{definition}[Lagrangien, problème dual]\label{def:qm:convex-optimisation:dual}
Pour la minimisation de $f(x)$ sous les contraintes $g_i(x) \le 0$ et $Ax = b$, le \emph{lagrangien}\index{lagrangien} est $L(x, \lambda, \nu) = f(x) + \sum_i\lambda_ig_i(x) + \nu^\top(Ax
- b)$, et la fonction duale est $d(\lambda, \nu) = \inf_xL(x, \lambda, \nu)$. Le \emph{problème dual}\index{problème dual} maximise $d$ sur $\lambda \ge 0$.
\end{definition}

\begin{proposition}[Dualité faible]\label{prop:qm:convex-optimisation:weak}
Pour tout $x$ réalisable et tout $\lambda \ge 0$, $d(\lambda, \nu) \le f(x)$~; par conséquent l'optimum dual est au plus égal à l'optimum primal.
\end{proposition}

\begin{proof}
Pour $x$ réalisable, $\lambda_ig_i(x) \le 0$ et $Ax - b = 0$, donc $L(x, \lambda, \nu) \le f(x)$, et $d(\lambda, \nu) \le L(x, \lambda, \nu)$.
\end{proof}

\begin{definition}[Dualité forte, condition de Slater, prix fictif]\label{def:qm:convex-optimisation:strong}
La \emph{dualité forte}\index{dualité forte} vaut lorsque les valeurs optimales primale et duale sont égales. La \emph{condition de Slater}\index{condition de Slater} est l'existence d'un
point strictement réalisable ($g_i(x) < 0$ pour les contraintes non linéaires, $Ax = b$). La variable duale optimale d'une contrainte est son \emph{prix fictif}\index{prix fictif}~: le
taux de variation de la valeur optimale lorsque la contrainte est relâchée.
\end{definition}

Pour un problème convexe, la \omterm{def:qm:convex-optimisation:strong}{condition de Slater} implique la \omterm{def:qm:convex-optimisation:strong}{dualité forte} et l'existence de variables duales optimales. Si la contrainte $g_i(x) \le 0$ est assouplie en $g_i(x) \le u$,
la valeur optimale $p^*(u)$ vérifie $\partial p^*/\partial u = -\lambda_i$ là où elle est dérivable~: un multiplicateur est la valeur marginale de la limite, exprimée dans l'unité de
l'objectif. Le relâchement mesure la même chose à l'échelle globale.

Sur le portefeuille du chapitre (\cref{fig:qm:convex-optimisation:turnover}), l'objectif est le rendement espéré moins cinq fois la variance (unités annuelles, covariance du
chapitre~22). Avec une limite de rotation de 20~\%, l'optimum présente un rendement espéré de 3.78~\%, une volatilité de 3.47~\% et une utilité de 3.18~\%~; la limite est saturée, et son multiplicateur vaut
0.0159~: chaque point de rotation supplémentaire autorisé vaut 1.6 point de base d'utilité. En résolvant à nouveau avec une limite de 10~\%, 15~\%, 30~\% et 40~\%, on trace la courbe dont la
pente est le multiplicateur~: 0.0210 à 10~\%, 0.0184 à 15~\%, 0.0119 à 30~\%, puis 0.0011 à 80~\% lorsque la limite cesse d'être contraignante.

\begin{omfigure}
\begin{tikzpicture}
\begin{axis}[width=10cm,height=5cm,xlabel={limite de rotation (\% du capital)},ylabel={\% annuel},
  tick label style={font=\small},label style={font=\small},xmin=0,xmax=85,ymin=2.6,ymax=4.4,grid=major,grid style={gray!25},
  legend columns=3,legend style={font=\footnotesize,at={(0.5,-0.3)},anchor=north,draw=none,fill=none,
  /tikz/every even column/.append style={column sep=6pt}},table/col sep=comma]
\addplot[omThm,very thick,mark=square*,mark size=1.5pt] table[x=limit,y=expret]{figdata/methods/23-convex-optimisation/turnover.csv};
\addplot[omDef,very thick,mark=*,mark size=1.5pt] table[x=limit,y=utility]{figdata/methods/23-convex-optimisation/turnover.csv};
\addplot[black,dashed,thick,domain=10:30,samples=2] {3.180 + 0.01593*(x-20)};
\legend{rendement espéré,utilité,pente à 20~\% = multiplicateur}
\end{axis}
\end{tikzpicture}

\omcaption{Le prix de la rotation~: l'utilité et le rendement espéré du portefeuille long-short optimal lorsque la limite de rotation est relâchée de 5~\% à 80~\% du capital. La ligne en pointillés
passant par le point à 20~\% a la pente du multiplicateur de la contrainte de rotation, 0.0159 (1.6 point de base d'utilité par 1~\% de rotation). Données~: le tutoriel du chapitre,
avec graine fixée.}\label{fig:qm:convex-optimisation:turnover}
\end{omfigure}

Relâcher chaque groupe de contraintes à tour de rôle et résoudre à nouveau les classe selon leur coût (\cref{fig:qm:convex-optimisation:costs})~: la limite de rotation 0.44~\% d'utilité, la limite
brute 0.28~\%, les quarante-trois limites de position saturées ensemble 0.19~\%, la neutralité sectorielle 0.014~\%, la neutralité en dollars rien. Ce dernier cas est instructif~: la neutralité en dollars est
impliquée par les cinq neutralités sectorielles, qui la somment, elle est donc redondante, et son multiplicateur (0.0004 dans la réponse du solveur) n'est pas unique~; une valeur de relâchement nulle avec un
multiplicateur non nul est la signature d'une contrainte redondante.

\begin{omfigure}
\begin{tikzpicture}
\begin{axis}[width=9cm,height=4.8cm,xbar,bar width=9pt,
  xlabel={utilité perdue à cause de la contrainte (\% annuel)},
  ytick={0,1,2,3,4},yticklabels={neutralité en dollars,neutralité sectorielle,limites de position,exposition brute,rotation},
  tick label style={font=\small},label style={font=\small},
  xmin=0,xmax=0.55,ymin=-0.6,ymax=4.6,grid=major,grid style={gray!25},
  xticklabel style={/pgf/number format/fixed,/pgf/number format/precision=2},scaled x ticks=false,table/col sep=comma]
\addplot[fill=omDef!60,draw=omDef,nodes near coords,nodes near coords style={font=\footnotesize,anchor=west,xshift=1pt,
  /pgf/number format/fixed,/pgf/number format/precision=3}]
  table[x=relax,y=k]{figdata/methods/23-convex-optimisation/costs.csv};
\end{axis}
\end{tikzpicture}

\omcaption{Ce que coûte chaque groupe de contraintes au portefeuille long-short~: le gain d'utilité lorsqu'il est retiré et que le problème est résolu à nouveau. La neutralité en dollars ne coûte rien, car les
neutralités sectorielles l'impliquent. Données~: le tutoriel du chapitre, avec graine fixée.}\label{fig:qm:convex-optimisation:costs}
\end{omfigure}

\begin{definition}[Certificat d'infaisabilité]\label{def:qm:convex-optimisation:certificate}
Un \emph{certificat d'infaisabilité}\index{certificat d'infaisabilité} pour le système linéaire $Ax = b$, $Gx \le h$ est un couple $(\lambda \ge 0, \nu)$ tel que $G^\top\lambda + A^\top\nu = 0$ et
$h^\top\lambda + b^\top\nu < 0$.
\end{definition}

Par le lemme de Farkas, un certificat existe si et seulement si le système est infaisable~: en combinant les lignes avec les poids $\lambda$ et $\nu$, on obtient $0 = (\lambda^\top G + \nu^\top A)x \le
h^\top\lambda + b^\top\nu < 0$, une impossibilité que chacun peut vérifier. Un problème de phase un, qui consiste à minimiser $t$ sous $Ax = b$, $Gx \le h + t\mathbf 1$, en produit un à partir de ses variables duales lorsque
son optimum est positif. Lorsqu'un responsable des risques exige que chacune des quarante actions du premier secteur soit détenue en position longue d'au moins 0.5~\%, alors que la neutralité sectorielle est activée, l'optimum de phase un
vaut $t^* = 0.005$, et le certificat affecte le poids $1/40$ à chacune des quarante lignes de position minimale ainsi qu'un multiplicateur à l'équation de neutralité du premier secteur, et rien à aucune
autre limite~: en additionnant les quarante lignes on obtient « le secteur est long d'au moins 20~\% », et la ligne de neutralité affirme qu'il l'est exactement à 0. Le rapport nomme les deux règles qui s'opposent.

\section{Conditions d'optimalité}

\begin{definition}[Conditions de Karush--Kuhn--Tucker]\label{def:qm:convex-optimisation:kkt}
Les \emph{conditions de Karush--Kuhn--Tucker}\index{conditions de Karush--Kuhn--Tucker} en $(x, \lambda, \nu)$ sont la faisabilité primale, la faisabilité duale ($\lambda \ge 0$), la complémentarité
($\lambda_ig_i(x) = 0$) et la stationnarité, $\nabla f(x) + \sum_i\lambda_i\nabla g_i(x) + A^\top\nu = 0$ (Karush, 1939~; Kuhn et Tucker, 1951).
\end{definition}

\begin{theorem}[Les conditions de KKT caractérisent l'optimum]\label{thm:qm:convex-optimisation:kkt}
Pour un problème convexe à fonctions différentiables vérifiant la \omterm{def:qm:convex-optimisation:strong}{condition de Slater}, $x$ est optimal si et seulement s'il existe $(\lambda, \nu)$ tel que $(x, \lambda, \nu)$ vérifie
les conditions de KKT~; les $(\lambda, \nu)$ sont alors des variables duales optimales.
\end{theorem}

\begin{proof}
Suffisance~: la stationnarité dit que $x$ minimise la \omterm{def:qm:convex-optimisation:convex}{fonction convexe} $L(\cdot, \lambda, \nu)$, donc $d(\lambda, \nu) = L(x, \lambda, \nu) = f(x)$ par complémentarité et
faisabilité~; la dualité faible fait alors de $x$ un optimum primal et de $(\lambda, \nu)$ un optimum dual. Nécessité~: sous Slater, la \omterm{def:qm:convex-optimisation:strong}{dualité forte} vaut avec un $(\lambda, \nu)$ optimal, et
$f(x) = d(\lambda, \nu) \le L(x, \lambda, \nu) \le f(x)$ force la complémentarité et la minimisation de $L$ en $x$, dont la condition du premier ordre est la stationnarité.
\end{proof}

La complémentarité a un sens économique~: une limite qui ne lie pas a un prix nul, et une limite dont le prix est positif est saturée. La réponse du solveur pour le portefeuille satisfait les
conditions à la précision de travail~: le résidu de stationnarité et la complémentarité sont tous deux inférieurs à $10^{-9}$. Les quarante-trois positions au niveau de la limite de 1.5~\% portent des multiplicateurs positifs,
d'environ 0.0175 chacun~; les 157 autres portent zéro.

\section{Programmes quadratiques et coniques du second ordre}

\begin{definition}[Programmes linéaires, quadratiques, coniques du second ordre et semi-définis]\label{def:qm:convex-optimisation:classes}
Un \emph{programme linéaire}\index{programme linéaire} minimise une fonction linéaire sous contraintes linéaires~; un \emph{programme quadratique}\index{programme quadratique} a un objectif quadratique convexe
et des contraintes linéaires~; un \emph{programme conique du second ordre}\index{programme conique du second ordre} (SOCP) ajoute des contraintes $\lVert A_ix + b_i\rVert_2 \le c_i^\top x + d_i$~; un
\emph{programme semi-défini}\index{programme semi-défini} impose à une matrice symétrique dépendant affinement de $x$ d'être semi-définie positive.
\end{definition}

Chaque classe contient la précédente. L'exposition brute et le taux de rotation, sommes de valeurs absolues, deviennent linéaires en décomposant chaque position en une partie positive et une partie négative~: le
portefeuille de 200 actions du chapitre est un \omterm{def:qm:convex-optimisation:classes}{programme quadratique} à 1\,000 variables.

\begin{proposition}[Risque et impact comme contraintes coniques]\label{prop:qm:convex-optimisation:cones}
Une limite de volatilité ou de tracking error $\sqrt{(w - b)^\top\Sigma(w - b)} \le \sigma_{\max}$ est la contrainte conique du second ordre $\lVert L^\top(w - b)\rVert_2 \le \sigma_{\max}$ avec $\Sigma =
LL^\top$. Un coût d'impact de marché croissant comme la puissance 3/2 du montant échangé, $|x|^{3/2} \le s$, se représente par deux contraintes coniques tournées~: il est vérifié si et seulement s'il existe
$r \ge 0$ avec $x^2 \le rs$ et $r^2 \le |x|$.
\end{proposition}

\begin{proof}
La première est l'identité $w^\top\Sigma w = \lVert L^\top w\rVert^2$. Pour la seconde (avec $x \ge 0$ après décomposition), si $x^{3/2} \le s$ on prend $r = \sqrt x$~: alors $r^2 = x$ et $x^2 = \sqrt x\,x^{3/2} \le rs$.
Réciproquement, $r^2 \le x$ donne $r \le \sqrt x$, donc $x^2 \le rs \le \sqrt x\,s$ et $x^{3/2} \le s$. Chaque contrainte $u^2 \le vw$ avec $v, w \ge 0$ est un cône du second ordre tourné,
$\lVert(2u, v - w)\rVert \le v + w$.
\end{proof}

Une limite de risque et une pénalité de risque sont le même problème vu de deux côtés. Maximiser le rendement espéré sous une limite de volatilité de 4~\% (neutre en dollars, positions limitées à 1.5~\%) avec le
solveur conique de la société, et maximiser le rendement moins $\gamma/2$ fois la variance, $\gamma$ étant trouvé par dichotomie pour donner une volatilité de 4~\%, produisent le même portefeuille (les poids diffèrent de
$3 \times 10^{-8}$), un rendement espéré de 5.31~\%, $\gamma = 25.8$, et un multiplicateur conique de 1.033 $= \gamma \times 0.04$~: le \omterm{def:qm:convex-optimisation:strong}{prix fictif} de la limite de volatilité, soit le rendement espéré par unité
de volatilité permise.

\begin{definition}[Matrice de corrélation la plus proche]\label{def:qm:convex-optimisation:ncm}
La \emph{matrice de corrélation la plus proche}\index{matrice de corrélation la plus proche} d'une matrice symétrique $C$ est la matrice semi-définie positive à diagonale unité la plus proche de $C$ au sens de la norme de Frobenius, un
\omterm{def:qm:convex-optimisation:classes}{programme semi-défini} que Higham (2002) résout par projections alternées avec la correction de Dykstra.
\end{definition}

Les corrélations estimées paire à paire à partir de données lacunaires ne sont pas garanties de former une matrice de corrélation, et une covariance construite à partir d'elles peut amener un optimiseur à promettre une variance négative.
Cinquante séries de soixante jours, dont la moitié des observations manquent, donnent une matrice paire à paire avec 21 valeurs propres négatives, la plus petite valant $-1.42$~; la \omterm{def:qm:convex-optimisation:ncm}{matrice de corrélation la plus proche}, atteinte en 74
projections, est à 4.45 au sens de la norme de Frobenius, ne modifie aucun coefficient de plus de 0.31, et n'a aucune valeur propre négative (\cref{fig:qm:convex-optimisation:ncm}).

\begin{omfigure}
\begin{tikzpicture}
\begin{axis}[width=10cm,height=4.8cm,xlabel={rang de la valeur propre},ylabel={valeur propre},
  tick label style={font=\small},label style={font=\small},xmin=0,xmax=51,ymin=-1.6,ymax=3,grid=major,grid style={gray!25},
  legend columns=2,legend style={font=\footnotesize,at={(0.5,-0.3)},anchor=north,fill=none,draw=none},table/col sep=comma]
\addplot[omThm,thick,mark=*,mark size=1.2pt] table[x=rank,y=before]{figdata/methods/23-convex-optimisation/ncm.csv};
\addplot[omDef,thick,mark=square*,mark size=1.2pt] table[x=rank,y=after]{figdata/methods/23-convex-optimisation/ncm.csv};
\addplot[black,thin,sharp plot] coordinates {(0,0) (51,0)};
\legend{corrélations paire à paire,matrice de corrélation la plus proche}
\end{axis}
\end{tikzpicture}

\omcaption{Valeurs propres (par ordre croissant~; la plus grande, supérieure à 10, hors du graphique) d'une matrice de corrélation paire à paire de cinquante séries dont la moitié des données manquent, et de sa
\omterm{def:qm:convex-optimisation:ncm}{matrice de corrélation la plus proche}. Les 21 valeurs propres négatives deviennent nulles. Données~: le tutoriel du chapitre, avec graine fixée.}\label{fig:qm:convex-optimisation:ncm}
\end{omfigure}

\section{Ce que font réellement les solveurs}

\begin{definition}[Méthode de points intérieurs]\label{def:qm:convex-optimisation:ipm}
Une \emph{méthode de points intérieurs}\index{méthode de points intérieurs} suit la trajectoire centrale~: elle applique la méthode de Newton aux conditions de KKT, la complémentarité étant relâchée en $s_iz_i = \mu$
pour les écarts $s$ et les multiplicateurs $z$, en gardant les deux strictement positifs, et fait tendre $\mu$ vers zéro.
\end{definition}

Chaque itération résout un seul système linéaire dans la direction de Newton~; le prédicteur-correcteur de Mehrotra (1992) calcule une direction affine, s'en sert pour choisir de combien réduire $\mu$, et corrige
le terme du second ordre, avec la même factorisation matricielle. Le nombre d'itérations dépend à peine de la taille du problème~: le portefeuille du chapitre converge en 18 itérations, les résidus
et la complémentarité passant d'un ordre de grandeur de un à moins de $10^{-9}$ (\cref{fig:qm:convex-optimisation:ipm}). Les \omterm{def:qm:convex-optimisation:ipm}{méthodes de points intérieurs} descendent de l'algorithme polynomial de Karmarkar (1984) pour
les \omterm{def:qm:convex-optimisation:classes}{programmes linéaires} et ont été étendues aux programmes coniques par Nesterov et Nemirovskii (1994). Les méthodes de découpage du premier ordre comme l'ADMM, dont le solveur OSQP (Stellato et coauteurs, 2020)
est un exemple de production, n'exigent qu'une seule factorisation et des itérations peu coûteuses, atteignent vite une précision modérée, et se réchauffent bien lorsque la solution de la veille est proche de celle du jour~; le solveur
conique de la société relève de cette famille, et a nécessité 924 itérations pour la limite de volatilité de 4~\% évoquée plus haut.

\begin{omfigure}
\begin{tikzpicture}
\begin{semilogyaxis}[width=10cm,height=5cm,xlabel={itération du point intérieur},ylabel={résidu},
  tick label style={font=\small},label style={font=\small},xmin=0,xmax=18,ymin=1e-10,ymax=100,grid=major,grid style={gray!25},
  legend columns=3,legend style={font=\footnotesize,at={(0.5,-0.3)},anchor=north,draw=none,fill=none,
  /tikz/every even column/.append style={column sep=8pt}},table/col sep=comma]
\addplot[omDef,very thick,mark=*,mark size=1.2pt] table[x=iter,y=dual]{figdata/methods/23-convex-optimisation/ipm.csv};
\addplot[omThm,very thick,mark=square*,mark size=1.2pt] table[x=iter,y=primal]{figdata/methods/23-convex-optimisation/ipm.csv};
\addplot[black,thick,dashed] table[x=iter,y=mu]{figdata/methods/23-convex-optimisation/ipm.csv};
\legend{stationnarité,faisabilité primale,complémentarité $\mu$}
\end{semilogyaxis}
\end{tikzpicture}

\omcaption{Convergence du solveur de points intérieurs de la société sur le portefeuille long-short (1\,000 variables)~: les normes des résidus de stationnarité et de faisabilité et la complémentarité moyenne
$\mu$ à chaque itération de Mehrotra, plafonnées à $10^{-9}$ (en dessous, les derniers chiffres dépendent de l'ordre des opérations en virgule flottante). Données~: le tutoriel du chapitre, avec graine fixée.}\label{fig:qm:convex-optimisation:ipm}
\end{omfigure}

\section{Tutoriel~: la contrainte qui a coûté le plus cher}\label{tut:qm:convex-optimisation:shadow}

\begin{tutorial}
\textbf{Objectif.} Formuler le portefeuille long-short comme un \omterm{def:qm:convex-optimisation:classes}{programme quadratique}, le résoudre avec la \omterm{def:qm:convex-optimisation:ipm}{méthode de points intérieurs} de la société, vérifier les conditions de KKT, lire les \omterm{def:qm:convex-optimisation:strong}{prix fictifs} et lire le
certificat d'un ensemble infaisable de limites.
\textbf{État final~:} \cref{fig:qm:convex-optimisation:turnover,fig:qm:convex-optimisation:costs,fig:qm:convex-optimisation:ipm}~; le multiplicateur du taux de rotation 0.0159 et le classement des
contraintes.
\begin{tutsteps}
\item \textbf{Le pas de Newton} de la \omterm{def:qm:convex-optimisation:ipm}{méthode de points intérieurs}, avec le prédicteur et le correcteur partageant une matrice.
\omcode{code/firm/portopt/firm_portopt.py}{50}{81}{L'itération de points intérieurs~: le prédicteur et le correcteur de Mehrotra partagent une matrice.}\label{lst:qm:convex-optimisation:ipm}
\item \textbf{Le problème de phase un} et son certificat de Farkas.
\omcode{code/firm/portopt/firm_portopt.py}{89}{107}{La phase un et le certificat d'infaisabilité.}\label{lst:qm:convex-optimisation:farkas}
\item \textbf{Exécuter} \texttt{today()}, \texttt{constraint\_costs()}, \texttt{turnover\_curve()}, \texttt{infeasible()}, \texttt{risk\_limit()} et \texttt{broken\_correlation()} dans
\texttt{qm\_opt.py}, puis \texttt{fig\_opt.py}.
\end{tutsteps}
\textbf{Pour aller plus loin.} Remplacer la limite sur le taux de rotation par un coût d'impact en puissance 3/2 dans l'objectif et comparer les montants échangés~; réchauffer la \omterm{def:qm:convex-optimisation:ipm}{méthode de points intérieurs} à partir de la solution de la veille
et compter les itérations économisées.
\end{tutorial}

\section{Construction~: l'optimiseur de portefeuille}\label{bld:qm:convex-optimisation:portopt}

\begin{build}
\bfield{Rôle} La construction de portefeuille de la mini-entreprise~: les prévisions, une \omterm{def:qm:covariance-estimation-and-random-matrices:cov}{matrice de covariance} (chapitre~22) et les limites en entrée, les ordres en sortie, avec un prix pour chaque limite et un certificat lorsque les limites
se contredisent.
\bfield{Interface} \texttt{qp(P, q, A, b, G, h)}~; \texttt{farkas(A, b, G, h)}~; \texttt{admm(P, q, C, lower, upper, cones)}~; \texttt{PortfolioProblem(alpha, Sigma, gamma, w0)} avec
\texttt{add\_budget}, \texttt{add\_neutral}, \texttt{add\_box}, \texttt{add\_gross}, \texttt{add\_turnover} et \texttt{solve()}~; \texttt{nearest\_correlation(C)}.
\bfield{Règles} Chaque résolution rapporte son statut, les résidus KKT et la variable duale de chaque contrainte nommée~; une demande non réalisable renvoie le certificat, jamais un portefeuille discrètement relâché~;
les covariances passent un contrôle de semi-définie positivité (et la correction par la \omterm{def:qm:convex-optimisation:ncm}{matrice de corrélation la plus proche}) avant utilisation.
\bfield{Tests d'acceptation} \texttt{code/firm/portopt/tests/}~: un petit QP et ses résidus KKT~; la forme fermée de la variance minimale et son multiplicateur~; les certificats de Farkas sur des systèmes réalisables et non réalisables~;
systèmes~; le solveur conique sur une boule de norme~; les contraintes du portefeuille sont vérifiées et le multiplicateur du taux de rotation égale la pente par différences finies~; la \omterm{def:qm:convex-optimisation:ncm}{matrice de corrélation la plus proche} est une matrice de corrélation et
laisse inchangée une matrice qui en est déjà une.
\bfield{Pour aller plus loin} Un plongement homogène auto-dual (des certificats sans phase un)~; la factorisation creuse du KKT pour des milliers d'actifs~; le coût d'impact en puissance 3/2 dans le solveur conique.
\end{build}

\begin{omsources}
\item W.~Karush, \emph{Minima of Functions of Several Variables with Inequalities as Side Constraints}, MSc thesis, University of Chicago, 1939~; H.~W.~Kuhn et A.~W.~Tucker, ``Nonlinear
programming'', \emph{Second Berkeley Symposium}, 1951.
\item H.~Markowitz, ``Portfolio selection'', \emph{Journal of Finance} 7, 1952.
\item N.~Karmarkar, ``A new polynomial-time algorithm for linear programming'', \emph{Combinatorica} 4, 1984~; S.~Mehrotra, ``On the implementation of a primal-dual interior point method'',
\emph{SIAM Journal on Optimization} 2, 1992.
\item Y.~Nesterov et A.~Nemirovskii, \emph{Interior-Point Polynomial Algorithms in Convex Programming}, SIAM, 1994.
\item S.~Boyd et L.~Vandenberghe, \emph{Convex Optimization}, Cambridge University Press, 2004.
\item M.~S.~Lobo, L.~Vandenberghe, S.~Boyd et H.~Lebret, ``Applications of second-order cone programming'', \emph{Linear Algebra and its Applications} 284, 1998.
\item N.~J.~Higham, ``Computing the nearest correlation matrix: a problem from finance'', \emph{IMA Journal of Numerical Analysis} 22, 2002.
\item B.~Stellato, G.~Banjac, P.~Goulart, A.~Bemporad et S.~Boyd, ``OSQP: an operator splitting solver for quadratic programs'', \emph{Mathematical Programming Computation} 12, 2020.
\end{omsources}

\section{Exercices}

\begin{exercise}[$\star$]\label{exo:qm:convex-optimisation:1}
La fonction $f(w) = \sqrt{w^\top\Sigma w}$ est-elle convexe lorsque $\Sigma$ est semi-définie positive~? Et $\{w: w^\top\Sigma w \ge 0.01\}$ est-il un \omterm{def:qm:convex-optimisation:convex}{ensemble convexe}~?
\end{exercise}

\begin{exercise}[$\star$]\label{exo:qm:convex-optimisation:2}
Le multiplicateur du taux de rotation vaut 0.0159. Estimer l'utilité gagnée en portant la limite de 20~\% à 22~\%, et comparer avec une nouvelle résolution (la courbe donne 3.180~\% à 20~\% et une pente de 0.0139 sur
[20~\%, 30~\%]).
\end{exercise}

\begin{exercise}[$\star$]\label{exo:qm:convex-optimisation:3}
Écrire $\sum_i|w_i - w_{0,i}| \le \tau$ avec des contraintes linéaires uniquement.
\end{exercise}

\begin{exercise}[$\star\star$]\label{exo:qm:convex-optimisation:4}
Déduire le \omterm{def:qm:covariance-estimation-and-random-matrices:pca}{portefeuille de variance minimale} et son multiplicateur des conditions KKT de $\min\frac12w^\top\Sigma w$ sous la contrainte $\mathbf 1^\top w = 1$.
\end{exercise}

\begin{exercise}[$\star\star$]\label{exo:qm:convex-optimisation:5}
Trouver un certificat de Farkas pour $x_1 + x_2 = 1$, $x_1 \le 0.2$, $x_2 \le 0.3$.
\end{exercise}

\begin{exercise}[$\star\star$]\label{exo:qm:convex-optimisation:6}
Montrer que $\max\alpha^\top w$ sous la contrainte $w^\top\Sigma w \le \sigma^2$ (aucune autre contrainte) admet pour solution $w = \sigma\Sigma^{-1}\alpha/\sqrt{\alpha^\top\Sigma^{-1}\alpha}$ et que le \omterm{def:qm:convex-optimisation:strong}{prix fictif} de la
limite de volatilité vaut $\sqrt{\alpha^\top\Sigma^{-1}\alpha}$.
\end{exercise}

\begin{exercise}[$\star\star\star$]\label{exo:qm:convex-optimisation:7}
\emph{Programmation.} Classer les groupes de contraintes du portefeuille du chapitre selon l'utilité que chacun coûte, en les relâchant un à un, et expliquer le résultat pour la neutralité en dollars.
\end{exercise}

\begin{exercise}[$\star\star\star$]\label{exo:qm:convex-optimisation:8}
\emph{Repérer l'erreur.} «~L'optimiseur a échoué ce soir, alors le travail de nuit a relâché les contraintes une à une jusqu'à trouver un portefeuille réalisable, et a échangé celui-ci.~»
\end{exercise}

\section{Problème~: la contrainte qui coûte le plus cher}

\begin{problem}[{Problème du week-end --- prix des limites d'un livre long-short}]\label{pb:qm:convex-optimisation:1}
Un livre de 200 actions (la covariance factorielle du chapitre~22, en unités annuelles) maximise le rendement espéré moins cinq fois la variance, neutre en dollars et par secteur, positions limitées à 1.5~\%, exposition brute
d'au plus 100~\%, taux de rotation d'au plus 20~\% par rapport au livre optimal de la veille. Les prévisions du jour sont celles de la veille amorties de 0.8, auxquelles s'ajoutent les nouvelles.

\medskip\noindent\textbf{Partie I --- L'optimum.}
\begin{enumerate}
\item Quels sont le rendement espéré, la volatilité et l'utilité du portefeuille optimal~?
\item Quelles contraintes sont saturées, et combien de positions se trouvent à leurs limites~?
\item Combien d'itérations de points intérieurs le solveur demande-t-il, et quels sont les résidus KKT~?
\item Quel est le multiplicateur de la limite de taux de rotation, et que signifie-t-il~?
\item Quel est le multiplicateur de l'exposition brute~?
\end{enumerate}

\medskip\noindent\textbf{Partie II --- Prix.}
\begin{enumerate}[resume]
\item Comment l'utilité évolue-t-elle lorsque la limite de taux de rotation passe de 5~\% à 80~\%~?
\item Comment les multiplicateurs se comparent-ils aux pentes de cette courbe~?
\item Classer les groupes de contraintes selon l'utilité que chacun coûte.
\item Pourquoi le coût de la neutralité en dollars est-il nul alors que son multiplicateur ne l'est pas~?
\item Quel rendement espéré est sacrifié pour chaque 1~\% de taux de rotation~?
\end{enumerate}

\medskip\noindent\textbf{Partie III --- Quand il n'existe aucun portefeuille.}
\begin{enumerate}[resume]
\item Que se passe-t-il lorsque chaque action du secteur~0 doit être détenue en position longue d'au moins 0.5~\%~?
\item Que dit le certificat, et comment le vérifie-t-on~?
\item Comment s'exprimerait une limite de volatilité de 4~\%, et quel est son \omterm{def:qm:convex-optimisation:strong}{prix fictif}~?
\item Pourquoi faut-il vérifier la matrice de corrélation avant d'optimiser~?
\item Que fait la \omterm{def:qm:convex-optimisation:ncm}{matrice de corrélation la plus proche} d'une matrice par paires comportant 21 valeurs propres négatives~?
\end{enumerate}

\medskip\noindent\textbf{Partie IV --- Jugement.}
\begin{enumerate}[resume]
\item Quelle limite faudrait-il demander au comité des risques d'assouplir en premier~?
\item Un prix de 1.6 point de base par 1~\% de taux de rotation vaut-il d'être payé, et par rapport à quoi~?
\item Que doit contenir chaque jour le rapport de l'optimiseur~?
\item Énoncer le \emph{résultat nommé}~: le \omterm{def:qm:convex-optimisation:strong}{prix fictif} de la limite de taux de rotation et le classement des contraintes selon le rendement qu'elles coûtent.
\item En une phrase~: qu'apporte la dualité à un gérant de portefeuille que la seule solution ne donne pas~?
\end{enumerate}
\end{problem}

\section{Questions d'entretien}

\begin{interviewq}[$\star$ \iqroles{chercheur, développeur}]\label{iq:qm:convex-optimisation:1}
Pourquoi la convexité est-elle si importante en optimisation~?
\end{interviewq}

\begin{interviewq}[$\star\star$ \iqroles{chercheur, trader}]\label{iq:qm:convex-optimisation:2}
Qu'est-ce qu'un \omterm{def:qm:convex-optimisation:strong}{prix fictif}, et comment l'utiliseriez-vous pour plaider en faveur d'un assouplissement d'une limite de risque~?
\end{interviewq}

\begin{interviewq}[$\star\star$ \iqroles{chercheur}]\label{iq:qm:convex-optimisation:3}
Énoncer les conditions KKT et expliquer la complémentarité en termes de portefeuille.
\end{interviewq}

\begin{interviewq}[$\star\star$ \iqroles{développeur, chercheur}]\label{iq:qm:convex-optimisation:4}
Votre optimiseur indique qu'un problème est non réalisable. Comment découvrez-vous quelles contraintes entrent en conflit~?
\end{interviewq}

\begin{interviewq}[$\star\star$ \iqroles{développeur}]\label{iq:qm:convex-optimisation:5}
\omterm{def:qm:convex-optimisation:ipm}{Méthodes de points intérieurs} ou méthodes du premier ordre pour un rééquilibrage quotidien de portefeuille de 3\,000 actions~: lesquelles, et pourquoi~?
\end{interviewq}

\begin{interviewq}[$\star\star\star$ \iqroles{chercheur}]\label{iq:qm:convex-optimisation:6}
Comment écrivez-vous une limite de tracking error et un coût d'impact en puissance 3/2 sous forme de contraintes convexes~?
\end{interviewq}
